El uso de IA ha sido un apoyo relevante para desarrollar este TFG dentro del tiempo disponible.
El volumen de código de la base proporcionada inicialmente por el ITC, que debía comprenderse,
sumado a la necesidad de trabajar con librerías y tecnologías nuevas y de cierta complejidad,
suponía un reto considerable para abarcar el proyecto en el plazo establecido.

La IA ha permitido acelerar el desarrollo, pero siempre bajo supervisión propia, validando y 
corrigiendo todo el contenido generado por la IA para garantizar su exactitud técnica y su autoría.
Durante el desarrollo del TFG se ha utilizado principalmente el modelo claude sonnet y el modelo local
gemma 4. Cabe destacar que la IA no ha sido utilizada para la toma de decisiones de diseño, 
integración y verificación del sistema, sino que ha sido utilizada como herramienta de apoyo para acelerar 
tareas repetitivas, explorar el código heredado, sugerir refactorizaciones y depurar, enfatizando esta última, 
ya que me ha permitido depurar errores en horas que de otra manera me hubieran llevado días, y en los peores casos, semanas.

En la \textbf{elaboración de esta memoria} se utilizó como ayuda de redacción
y estructuración a partir del código y la documentación reales del proyecto, permitiendo mejorar mucho
la calidad, especialmente de los gráficos y tablas. Asimismo, ha sido de gran utilidad para la recopilación
de información y referencias bibliográficas, así como para la generación de resúmenes 
y explicaciones de conceptos complejos.

Finalmente, el proyecto incorpora un asistente LLM preparado para funcionar con gemma 4, claude, openai y compatibles.
